\sectioncentered*{Введение}
\addcontentsline{toc}{section}{Введение}

Современный этап развития мировых социально-экономических отношений характеризуется беспрецедентным уровнем глобализации и устойчивым ростом интенсивности международных перемещений. Несмотря на периодические колебания, вызванные глобальными кризисами, общая тенденция к расширению географии деловых поездок и индивидуального туризма сохраняет свою актуальность. В этих условиях обеспечение безопасности путешественников становится одной из приоритетных задач как для корпоративного сектора, так и для специализированных туристических компаний. Развитие концепции «должной осмотрительности» (Duty of Care), закрепленной в международных стандартах, таких как ISO 31030, возлагает на организации юридическую и моральную ответственность за жизнь и здоровье своих клиентов и сотрудников во время их нахождения за рубежом.

Ландшафт глобальных угроз в последние годы претерпел существенные изменения, став значительно более динамичным, многофакторным и труднопредсказуемым. Путешественники ежедневно сталкиваются с широким спектром рисков: от внезапных геополитических конфликтов, массовых протестов и террористических актов до стремительно развивающихся природных катаклизмов, вызванных климатическими изменениями. Эпидемиологические угрозы, сбои в работе критической транспортной инфраструктуры и возрастающий уровень локальной криминогенной напряженности требуют непрерывного контроля. В такой нестабильной среде традиционные методы оценки рисков, опирающиеся на статические страновые отчеты и периодические сводки, демонстрируют свою полную несостоятельность.

Анализ достижений в области управления рисками в путешествиях показывает, что существующие на рынке решения исторически развивались по двум основным направлениям. Первое направление базируется на экспертно-аналитическом подходе, при котором штат специалистов вручную осуществляет мониторинг новостных лент, верификацию информации и рассылку предупреждений. Данный метод обеспечивает высокую точность и отсутствие ложных срабатываний, однако характеризуется критически низкой скоростью реакции (задержки могут достигать нескольких часов) и колоссальными операционными издержками, что делает такие системы недоступными для малого бизнеса и индивидуальных туристов. Второе направление использует алгоритмический сбор данных на основе жестко заданных ключевых слов. Несмотря на высокую скорость обработки информации, подобные агрегаторы генерируют неприемлемый уровень информационного шума, поскольку алгоритмы не способны улавливать контекст, тональность сообщения и отличать реальную угрозу от исторической справки или метафоры в тексте новостной статьи.

Настоящим прорывом в решении проблемы глобального мониторинга стало стремительное развитие технологий искусственного интеллекта, в частности, методов обработки естественного языка (Natural Language Processing, NLP) и больших языковых моделей (Large Language Models, LLM). Появление архитектур на базе трансформеров позволило перейти от простого синтаксического поиска к глубокому семантическому анализу неструктурированных текстовых данных. Использование методов Zero-shot классификации дает возможность сопоставлять поступающие новостные сообщения с обширной таксономией угроз без необходимости постоянного переобучения моделей на новых наборах данных. Интеграция алгоритмов продвинутого геопарсинга позволяет с высокой точностью извлекать из текста пространственные координаты описываемых событий, формируя зоны потенциального поражения. Совокупность этих инновационных подходов открыла возможность создания систем мониторинга, способных обрабатывать огромные массивы информации из систем открытой разведки (OSINT), таких как проект GDELT, в режиме реального времени, извлекая только релевантные данные и сводя к минимуму участие человека в процессе первичной фильтрации.

Целью дипломного проектирования является разработка интеллектуальной системы мониторинга рисков и обеспечения безопасности путешественников по индивидуальному заказу туристической компании ООО «ТревелСейф».

Реализация данной цели направлена на автоматизацию процесса сбора, семантического анализа и оперативного оповещения об угрозах для обеспечения принципа «должной осмотрительности» и надежной защиты клиентов компании во время их поездок.

В основу проектирования, научного исследования и поиска технического решения положены следующие принципы:
\begin{itemize}
\item использование событийно-ориентированной микросервисной архитектуры для обеспечения высокой отказоустойчивости и сверхвысокой скорости обработки независимых потоков данных;
\item применение передовых моделей машинного обучения для глубокого семантического анализа открытых источников без необходимости ручной модерации входящего контента;
\item внедрение алгоритмов точного геопарсинга для автоматического извлечения пространственных координат инцидентов из неструктурированных текстовых массивов;
\item интеграция механизмов геофенсинга для непрерывного сопоставления текущего местоположения пользователя или его планируемого маршрута с динамически изменяющимися зонами риска;
\item применение локальных больших языковых моделей для автоматического реферирования новостных сводок и генерации лаконичных уведомлений;
\item изоляция промежуточных сырых данных от очищенной аналитической информации на уровне структуры базы данных для обеспечения чистоты анализа;
\item ориентация на конечного пользователя посредством создания интуитивно понятного интерфейса с интерактивной глобальной картой и гибкой системой настройки профиля угроз.
\end{itemize}

Для достижения поставленной цели в рамках дипломного проектирования необходимо решить следующие задачи:
\begin{itemize}
\item провести комплексный анализ предметной области и существующих на мировом рынке систем управления рисками в путешествиях;
\item выявить недостатки текущих аналогов и обосновать необходимость применения интеллектуального подхода на базе нейронных сетей;
\item выполнить сравнительный анализ источников данных о глобальных инцидентах и выбрать наиболее оптимальные агрегаторы;
\item обосновать выбор архитектурного паттерна, технологического стека и конкретных моделей машинного обучения для реализации вычислительного ядра системы;
\item спроектировать микросервисную архитектуру программного комплекса и разработать схемы взаимодействия его внутренних компонентов;
\item разработать структуру реляционной базы данных с поддержкой пространственных запросов для хранения информации об инцидентах и пользователях;
\item рассчитать экономические показатели целесообразности инвестиций в разработку и внедрение программного средства для организации-заказчика.
\end{itemize}

Пояснительная записка к дипломному проекту структурно состоит из введения, трех основных разделов и заключения.

Краткое изложение содержания разделов пояснительной записки:
\begin{itemize}
\item в первом разделе проводится анализ подходов к разработке интеллектуальной системы мониторинга рисков, исследуются существующие на рынке аналогичные решения, подробно рассматриваются источники данных о глобальных событиях, а также обосновывается выбор технологий, фреймворков и моделей машинного обучения;
\item во втором разделе приведено описание функциональных требований к системе, выполнено архитектурное проектирование ее микросервисных компонентов, а также представлено детальное описание структуры базы данных, программных интерфейсов и алгоритмов взаимодействия модулей сбора, анализа и оповещения;
\item в третьем разделе осуществляется экономическое обоснование целесообразности инвестиций в разработку программного средства, включающее расчет затрат на заработную плату разработчиков, формирование отпускной цены системы, а также расчет показателей экономической эффективности, таких как чистый дисконтированный доход и срок окупаемости;
\item в заключении подводятся итоги проделанной работы, формулируются основные выводы по результатам проектирования и экономического обоснования, а также подтверждается успешное достижение поставленной цели дипломного проекта.
\end{itemize}

Дипломный проект выполнен самостоятельно, проверен в системе "Антиплагиат". Процент оригинальности соответствует норме, установленной кафедрой ИИТ. Цитирования обозначены ссылками на публикации, указанные в "Списке использованных источников".